\section{Exactness under Monotone Contracts}
\label{sec:monotone}

\(\PsiDeletionMonotonicity\) says that deleting more contract atoms cannot
destroy reference feasibility.  This is natural for standard conjunctive
axiom-deletion contracts, where deleting an atom removes constraints, but it is
not assumed for arbitrary reportability contracts.

\begin{theorem}[M3-C: monotone converse]
\label{thm:m3c}
If \(\PsiDeletionMonotonicity\) holds and grouped correctness holds pointwise,
\[
  \forall G,\quad
  \GroupedRepair(G)\leftrightarrow\ContractRepair(G),
\]
then \(\GroupSoundness\) holds.
\end{theorem}

The Lean declaration \code{m3c_converse} does not assume residual faithfulness.
It descends from a raw repair's grouped image to a minimal grouped report, uses
grouped correctness to obtain a contract repair, and then applies reference
deletion monotonicity.

\begin{corollary}[Exactness]
\label{cor:exactness}
For a residually faithful realization with pairwise disjoint active blocks over
a \(\Psi\)-deletion-monotone contract,
\[
  \GroupSoundness
  \quad\Longleftrightarrow\quad
  \forall G,\
  \GroupedRepair(G)\leftrightarrow\ContractRepair(G).
\]
\end{corollary}

This is \code{groupSoundness_iff}.  Residual faithfulness is used in the M3-B
forward direction, not in the M3-C converse.

\begin{proposition}[Audit-cost collapse]
\label{prop:audit-collapse}
Under \(\PsiDeletionMonotonicity\), if every raw minimal repair \(R\) satisfies
\[
  \SatPsi(I\setminus\groupTouchAny(R)),
\]
then unrestricted \(\GroupSoundness\) holds.
\end{proposition}

The Lean declaration is \code{audit_cost_collapse}.  Its finite descent does
not assume deletion monotonicity of \(\SatPhi\).

The monotonicity assumption is not removable.  The compiled finite example in
\code{Examples.lean} has pointwise grouped correctness but fails both
\(\GroupSoundness\) and \(\PsiDeletionMonotonicity\).

\paragraph{Comparison.}
Table~\ref{tab:theorem-comparison} separates the roles of the three theorem
components.

\begin{table}[t]
\centering
\scriptsize
\begin{tabularx}{\textwidth}{@{}p{0.07\textwidth}p{0.20\textwidth}p{0.20\textwidth}p{0.27\textwidth}p{0.08\textwidth}X@{}}
\toprule
Component & Interface / soundness condition & Other assumptions & Conclusion & Raw transport? & Lean status \\
\midrule
M3-A & \(\RepairAtomicity\) & block disjointness; \(\ResidualFaithfulness\) & pointwise grouped correctness & yes, atom-indexed & checked \\
M3-B & \(\GroupSoundness\) & block disjointness; \(\ResidualFaithfulness\) & pointwise grouped correctness; two-realization corollary & no & checked \\
M3-C & pointwise grouped correctness & \(\PsiDeletionMonotonicity\) & \(\GroupSoundness\) & no & checked \\
\bottomrule
\end{tabularx}
\caption{M3-A, M3-B, and M3-C have different logical roles.}
\label{tab:theorem-comparison}
\end{table}
